% Generated draft: requires mathematical and source review.
\documentclass[11pt]{article}
\usepackage[margin=25mm]{geometry}
\usepackage{fontspec}
\setmainfont{Latin Modern Roman}
\usepackage{amsmath,amssymb,mathtools,unicode-math}
\setmathfont{Latin Modern Math}
\usepackage{enumitem,needspace,xurl,hyperref}
\setlength{\parindent}{0pt}
\setlength{\parskip}{5pt}
\setlength{\emergencystretch}{3em}
\newcommand{\N}{\mathbb{N}}
\newcommand{\Z}{\mathbb{Z}}
\newcommand{\Q}{\mathbb{Q}}
\newcommand{\R}{\mathbb{R}}
\newcommand{\C}{\mathbb{C}}
\newcommand{\F}{\mathbb{F}}
\newcommand{\A}{\mathbb{A}}
\newcommand{\PP}{\mathbb{P}}
\newcommand{\E}{\mathbb{E}}
\newcommand{\eps}{\varepsilon}
\newcommand{\dd}{\,\mathrm{d}}
\begin{document}
\section*{UnsolvedMath ID: 30004494}
\textbf{Source problem code:} catalogue entry
\subsection{Definitions and mathematical statement}
Let \(B\) be smooth and projective, and let \(Z=\bigcup_{i=1}^{r}Z_i\) be a simple normal-crossing divisor, with irreducible components \(Z_i\). Thus the components are smooth and meet transversely; locally the divisor is given by a product of distinct coordinate functions. Put \(U=B\setminus Z\).

In the period-map setting intended by the supplied record, let \(\Lambda\) denote the Hodge line bundle on \(U\), and let \(\Lambda^e\) denote its canonical extension to \(B\). Assume, as in the supplied clean statement, that the period map has injective logarithmic differential. These are source-dependent hypotheses: the record does not specify the base field, the variation of Hodge structure or period domain, the construction of the Hodge line bundle, the canonical-extension convention, or the precise meaning of injectivity of the logarithmic differential. The question below is conditional on these data being specified so that \(\Lambda^e\) is a line bundle on \(B\).

For integers \(m,a_1,\ldots,a_r\), interpret the additive line-bundle notation by
\[
 m\Lambda^e-\sum_{i=1}^{r}a_i[Z_i]
 \;:=\;
 (\Lambda^e)^{\otimes m}\otimes
 \mathcal O_B\!\left(-\sum_{i=1}^{r}a_iZ_i\right).
\]
Here ampleness means that some positive tensor power of the line bundle is very ample.

\textbf{Supplied ampleness question.} For every choice of the above data satisfying the intended hypotheses, do there exist fixed integers \(a_1,\ldots,a_r\geq 0\) and \(m_0\in\Z\) such that
\[
 \forall m\in\Z,\qquad
 m\geq m_0\ \Longrightarrow\
 (\Lambda^e)^{\otimes m}\otimes
 \mathcal O_B\!\left(-\sum_{i=1}^{r}a_iZ_i\right)
 \text{ is ample on }B?
\]
The coefficients \(a_i\) and the threshold \(m_0\) may depend on the given geometric and period-map data, but the coefficients must be independent of \(m\).

\textbf{Unresolved qualification.} The original formulation assumes ``suitable local-Torelli-type assumptions,'' whereas the clean formulation specifies injectivity of the logarithmic differential. The supplied record does not give the additional hypotheses or establish the equivalence of these formulations. Consequently, the displayed question records the supplied target without claiming that its full hypothesis set has been reconstructed.
\subsection{Short English statement}
For the supplied period-map data on \(B\setminus Z\), can one choose a fixed effective integral divisor supported on the boundary so that every sufficiently high tensor power of the extended Hodge line bundle, twisted by the negative of that divisor, is ample? The exact period-map and local-Torelli hypotheses require clarification. The supplied partially solved label is dated curation, not an independently verified conclusion.
\subsection{Sources}
\begin{itemize}
\item Local UnsolvedMath registry, permanent ID \(30004494\), title \emph{Boundary-Corrected Ampleness of Extended Hodge Bundles}; source problem code \texttt{catalogue entry}, in the Oberwolfach Reports Open Problems collection. The supplied proposed year is \(2020\); no individual proposer is named.
\item Supplied source citation: \emph{Komplexe Analysis - Algebraicity and Transcendence} (2021), Oberwolfach Reports, DOI \texttt{10.4171/owr/2020/24}, pp.\ 1263--1313. \url{https://doi.org/10.4171/owr/2020/24}. This is also listed among the supplied literature sources as \emph{Oberwolfach Report 24/2020}. The citation's year and report designation are retained as supplied.
\item Supplied literature source: \emph{Natural line bundles on completions of period mappings}. \url{https://arxiv.org/abs/2102.06310}. The record supplies no authors or theorem locations for this reference.
\item Literature assessment dated exactly \texttt{2026-08-22}: the supplied curation labels the problem partially solved and states that the ampleness conjecture is proved when the boundary is irreducible and in dimension two under stronger hypotheses, while the general normal-crossing-boundary assertion remains unresolved. The stronger hypotheses and precise theorem statements are not supplied. This is supplied dated curation, not independently verified; the record itself requests independent verification of time-sensitive status and cited claims.
\item The registry labels the statement \texttt{corrected\_verified} and reports that it was checked against the official Oberwolfach report. This verification claim is part of the supplied record; no independent source consultation or mathematical verification is asserted here.
\end{itemize}
\end{document}
